\begin{acknowledgments}

The author acknowledges the use of ChatGPT (OpenAI) as an assistive
tool during software development, code debugging, numerical workflow
development, and manuscript preparation. All scientific results,
interpretations, and final manuscript content were reviewed and
verified by the author.

% Eric will add funding information and institutional support here.

\end{acknowledgments}

\section*{AUTHOR DECLARATIONS}

\subsection*{Conflict of Interest}

The authors have no conflicts to disclose.

\subsection*{Author Contributions}

A.C.-G. developed the finite-pulse Liouvillian dressing formalism,
implemented the PULSUS software, performed the numerical calculations
and validation studies, analyzed the results, prepared the figures,
and wrote the original draft of the manuscript. E.R.B. proposed the
initial use of Liouvillian Green's-function methods for the
frequency-domain formulation, supervised the project, and contributed
to the conceptual development, interpretation of the results, and
revision of the manuscript. Both authors discussed the results and
approved the final version of the manuscript.

\section*{DATA AVAILABILITY}

PULSUS is open-source software distributed under the BSD 3-Clause
license. The source code is publicly available through GitHub, and the
package can be installed from the Python Package Index as
\texttt{pulsus-spec}. The version used for the calculations reported
in this work, PULSUS v0.2.0, is archived on Zenodo under the
version-specific DOI 10.5281/zenodo.23253040. Archived releases are
also available through the Zenodo concept DOI
10.5281/zenodo.22755139.

